\subsection{Universally measurable functional calculus}

In this section the Hilbert spaces considered are complex.

Let u be a partial normal operator on a Hilbert space E. We denote by \(\mu_{x,y}\) (resp. \(\mu_{y}\) ) the spectral measure of \((x,y)\in\mathrm{E}\times\mathrm{E}\) (resp. of y) relative to u.

Let \(y \in \mathrm{E}\) . The map \(\mathcal{K}(\mathrm{Sp}(u)) \to \mathrm{E}\) defined by \(f \mapsto f(u)(y)\) satisfies

\[
\| f (u) (y) \| ^ {2} = \int_ {\mathrm{Sp} (u)} | f | ^ {2} \mu_ {y} = \| f \| _ {\mathcal {L} ^ {2} (\mathrm{Sp} (u), \mu_ {y})} ^ {2}.
\]

There therefore exists a unique isometric linear map \(ev_{y}\) from the space \(\mathrm{L}^{2}(\mathrm{Sp}(u),\mu_{y})\) into E such that \(\mathrm{ev}_{y}(\widetilde{f})=f(u)(y)\) if \(\widetilde{f}\) is the class of a function \(f\in\mathcal{K}(\mathrm{Sp}(u))\) .

Let \(f\) be a universally measurable function defined on the spectrum of \(u\) . We denote by \(\mathrm{D}_{f}\) the set of elements \(y \in \mathrm{E}\) such that \(f\) belongs to \(\mathcal{L}^{2}(\mathrm{Sp}(u), \mu_{y})\) .

PROPOSITION 3. --- Let f be a universally measurable function defined on the spectrum of u. The set D \(_{f}\) is a dense vector subspace of E. The map y ↦ ev \(_{y}\) (f) is a partial normal operator on E whose domain is D \(_{f}\) , and which is denoted f(u).

For every \(x \in \mathrm{E}\) and every \(y \in \mathrm{D}_f\) , we have \(f \in \mathcal{L}^1(\mathrm{Sp}(u), \mu_{x,y})\) and

\[
\langle x \mid f (u) y \rangle = \int_ {\operatorname{Sp} (u)} f \mu_ {x, y}.\tag{6}
\]

We shall prove the following more precise statement:

PROPOSITION 4. --- Let f be a universally measurable function on the spectrum of u.

\begin{enumerate}
\def\labelenumi{\alph{enumi})}
\item
  The set \(D_{f}\) is a dense vector subspace of E. The map \(f(u): y \mapsto \operatorname{ev}_{y}(f)\) from \(D_{f}\) into E is linear and coincides with \(1_{E}\) if f = 1;
\item
  Let \(g \in \mathcal{L}_{\mathrm{u}}(\mathrm{Sp}(u))\) . For every \(y \in \mathrm{D}_f\) and every \(x \in \mathrm{D}_g\) , we have \(fg \in \mathcal{L}^1(\mathrm{Sp}(u), \mu_{x,y})\) , and
\end{enumerate}

\[
\langle g (u) x \mid f (u) y \rangle = \int_ {\mathrm{Sp} (u)} \overline {{g}} f \mu_ {x, y}.\tag{7}
\]

\begin{enumerate}
\def\labelenumi{\alph{enumi})}
\setcounter{enumi}{2}
\tightlist
\item
  Suppose that \(E = L^{2}(X, \mu)\) where X is a locally compact topological space and \(\mu\) is a positive measure on X, and that \(u = m_{h}\) , where \(h: X \to \mathbf C\) is \(\mu\) -measurable. We have \(f(m_{h}) = m_{f \circ h}\) .
\end{enumerate}

By Theorem 1 of IV, p. 266, we may assume that we are in the situation of assertion \(c\) ), namely, \(\mathrm{E} = \mathrm{L}^2 (\mathrm{X},\mu)\) and \(u = m_h\) , where X is a locally compact topological space, \(\mu\) is a positive measure on X and \(h\colon \mathrm{X}\to \mathbf{C}\) is \(\mu\) -measurable. We shall denote by \(\widetilde{\varphi}\) the class in \(\mathrm{L}^2 (\mathrm{X},\mu)\) of a function \(\varphi \in \mathcal{L}^2 (\mathrm{X},\mu)\) .

Let \(\mathrm{S} = \mathrm{Sp}(m_h)\) . We denote by \(\mu_{\widetilde{\varphi}_1, \widetilde{\varphi}_2}\) (resp. \(\mu_{\widetilde{\varphi}}\) ) the spectral measure of \((\widetilde{\varphi}_1, \widetilde{\varphi}_2)\) (resp. of \(\widetilde{\varphi}\) ) relative to \(m_h\) for all \((\varphi_1, \varphi_2) \in \mathcal{L}^2(\mathrm{X}, \mu)^2\) (resp. every \(\varphi \in \mathcal{L}^2(\mathrm{X}, \mu)\) ).

Let \(\varphi\in\mathcal{L}^{2}(X,\mu)\) . The spectral measure \(\mu_{\widetilde{\varphi}}\) is equal to the image measure \(h(|\varphi|^{2}\cdot\mu)\) on S (Lemma 4 of IV, p. 269). Since the function \(f\) is \(\mu_{\varphi}\) -measurable, we have \(\varphi\in D_{f}\) if and only if the integrals

\[
\int_ {\mathrm{S}} ^ {*} | f | ^ {2} d \mu_ {\widetilde {\varphi}} = \int_ {\mathrm{S}} ^ {*} | f | ^ {2} h (| \varphi | ^ {2} d \mu) = \int_ {\mathrm{X}} ^ {*} | f \circ h | ^ {2} | \varphi | ^ {2} d \mu
\]

are finite (INT, V, p. 70, § 6, no. 2, prop. 2). This means that \(D_{f}\) is the domain of the multiplication operator \(m_{f\circ h}\) , which is a dense subspace of E (prop. 5 of IV, p. 232).

The restriction of \(\mathrm{ev}_{\widetilde{\varphi}}\) to the function classes \(g \in \mathcal{K}(\mathrm{S})\) is the map which associates to a class \(\widetilde{g}\) the element \(g(m_h)(\widetilde{\varphi}) = m_{g \circ h}(\widetilde{\varphi})\) (Lemma 4 of IV, p. 269). The map \(\mathrm{ev}_{\widetilde{\varphi}}\) therefore coincides with the isometric map from \(\mathrm{L}^2(\mathrm{S}, \mu_{\widetilde{\varphi}})\) into \(\mathrm{L}^2(\mathrm{X}, \mu)\) obtained by passing to quotients from the map \(g \mapsto (g \circ h) \cdot \varphi\) from \(\mathcal{L}^2(\mathrm{S}, \mu_{\widetilde{\varphi}})\) into \(\mathcal{L}^2(\mathrm{X}, \mu)\) .

In particular, we therefore have \(\mathrm{ev}_{\widetilde{\varphi}}(f)=m_{f\circ h}(\widetilde{\varphi})\) , so the map \(f(m_{h})\) from \(D_{f}\) into E coincides with the partial operator \(m_{f\circ h}\) . This proves assertion c); if f=1, we find \(\mathrm{ev}_{\widetilde{\varphi}}(1)=\widetilde{\varphi}\) , which also completes the proof of a).

Let us prove assertion \(b\) ). Let \(\varphi_{1}\) and \(\varphi_{2}\) be functions such that \(\widetilde{\varphi}_{1} \in \mathrm{D}_{f}\) and \(\widetilde{\varphi}_{2} \in \mathrm{D}_{g}\) . It follows from (INT, V, loc. cit.)

\[
\begin{array}{l} \int_ {\mathrm{S}} ^ {*} | f g |   | \mu_ {\widetilde {\varphi_ {1}}, \widetilde {\varphi_ {2}}} | = \int_ {\mathrm{X}} ^ {*} | (f \circ h) (g \circ h) \varphi_ {1} \varphi_ {2} | d \mu \\ \qquad \leqslant \| (f \circ h) \varphi_ {1} \| \| (g \circ h) \varphi_ {2} \| \end{array}
\]

which is finite, hence \(fg \in \mathcal{L}^1(\mathrm{S}, \mu_{\widetilde{\varphi}_1, \widetilde{\varphi}_2})\) . Moreover, we then have

\[
\begin{array}{c} \langle g (m _ {h}) (\widetilde {\varphi} _ {2})   |   f (m _ {h}) (\widetilde {\varphi} _ {1}) \rangle = \int_ {\mathrm{X}} \overline {{(g \circ h)   \varphi_ {2}}}    (f \circ h)   \varphi_ {1}    d \mu \\ = \int_ {\mathrm{S}} \overline {{g}}   f    d \mu_ {\widetilde {\varphi} _ {1}, \widetilde {\varphi} _ {2}}, \end{array}
\]

which concludes the proof.

DEFINITION 5. --- The map \(f \mapsto f(u)\) from \(\mathcal{L}_{\mathrm{u}}(\mathrm{Sp}(u))\) into the set of normal partial operators on E defined by prop. 3 of IV, p. 271 is called the universally measurable functional calculus mapping associated with \(u\) .

More generally, let T be a set and let \(g\colon\operatorname{Sp}(u)\times\operatorname{T}\to\mathbf{C}\) be a map; let \(t\in T\) be such that the function \(g_{t}\colon z\mapsto g(z,t)\) is universally measurable on \(\operatorname{Sp}(u)\) . We then denote by \(g(u,t)=g_{t}(u)\) .

Formula (7) implies in particular that

\[
\| f (u) (y) \| ^ {2} = \int_ {\operatorname{Sp} (u)} | f | ^ {2} \mu_ {y}\tag{8}
\]

for all \(f \in \mathcal{L}_{\mathrm{u}}(\mathrm{Sp}(u))\) and every \(y \in \mathrm{D}_f\) . From this, taking \(f = 1\) we deduce that \(\mu_y\) is a positive measure of total mass \(\| y\|^2\) .

If \(u\) is a normal partial operator on a Hilbert space E and \(f \in \mathcal{K}(\mathrm{Sp}(u))\) , then \(\mathrm{D}_f = \mathrm{E}\) and the partial operator \(f(u)\) is then continuous, since

\[
\| f (u) y \| \leqslant \| f \| _ {\infty} \| y \|
\]

for all \(y \in E\) according to (8). The endomorphism \(f(u)\) coincides with that of Definition 3 of IV, p. 268.

COROLLARY 1. --- a) For every \(f \in \mathcal{L}_{\mathrm{u}}(\mathrm{Sp}(u))\) , the partial operator \(f(u)\) is normal and \(f(u)^{*} = \overline{f}(u)\) . Moreover \(f(u)\) is positive if \(f \geqslant 0\) , and self-adjoint if \(f\) takes real values;

\begin{enumerate}
\def\labelenumi{\alph{enumi})}
\setcounter{enumi}{1}
\item
  Let \(k \in \mathbf{N}\) and \(f(z) = z^{k}\) for \(z \in \operatorname{Sp}(u)\) . We have \(f(u) = u^{k}\) ;
\item
  Let \(\lambda \in \mathbf{C} - \mathrm{Sp}(u)\) and \(f(z) = (\lambda - z)^{-1}\) for \(z \in \mathrm{Sp}(u)\) . We have \(f(u) = \mathrm{R}(u, \lambda)\) ;
\item
  One has \(\beta(u) = b(u)\) ;
\item
  Let \(f \in \mathcal{L}_{\mathrm{u}}(\mathrm{Sp}(u))\) and \(g \in \mathcal{L}_{\mathrm{u}}(\mathrm{Sp}(u))\) such that \(|g| \leqslant 1 + |f|\) . The domain of \(g(u)\) is a core for \(f(u)\) .
\end{enumerate}

In view of the spectral theorem (th. 1 of IV, p. 266), this follows from the preceding proposition combined, respectively, with:

\begin{enumerate}
\def\labelenumi{\alph{enumi})}
\item
  lemma 2, a) of IV, p. 266, prop. 23 of IV, p. 253, and its corollary;
\item
  prop. 24, b) of IV, p. 255;
\item
  Prop. 22, b) of IV, p. 252;
\item
  Lemma 2, a) of IV, p. 266;
\item
  Prop. 6, b) of IV, p. 232.
\end{enumerate}

Note. --- For \(f = \operatorname{Id}_{\operatorname{Sp}(u)}\) , we have \(f(u) = u\) (assertion b) for \(k = 1\) ). The domain of u therefore coincides with the set of \(x \in E\) such that the identity map of \(\operatorname{Sp}(u)\) belongs to \(\mathcal{L}^{2}(\operatorname{Sp}(u), \mu_{x})\) ; it contains in particular the elements \(x \in E\) such that the measure \(\mu_{x}\) is compactly supported.

The following corollary generalizes the corollary of Prop. 16 of IV, p. 247 and Prop. 17 of IV, p. 248.

COROLLARY 2. --- Let u be a partial normal operator on the Hilbert space E. Assume that E is nonzero.

\begin{enumerate}
\def\labelenumi{\alph{enumi})}
\item
  The spectrum of u is nonempty;
\item
  For every \(\lambda \in \mathbf{C} - \mathrm{Sp}(u)\) , the norm of the resolvent \(\mathrm{R}(u, \lambda)\) is equal to \(1/\delta\) , where \(\delta > 0\) is the distance in \(\mathbf{C}\) from \(\lambda\) to the spectrum of \(u\) .
\item
  For every \(\varepsilon > 0\) , the \(\varepsilon\) -pseudo-spectrum \(\mathrm{PSp}_{\varepsilon}(u)\) is the set of \(\lambda \in \mathbf{C}\) at distance \(< \varepsilon\) from \(\mathrm{Sp}(u)\) .
\end{enumerate}

Suppose that the spectrum of \(u\) is empty. Then \(u\) is injective, and the endomorphism \(u^{-1} = -\mathrm{R}(u,0)\) is normal (cor. 1, c) and a)). Since \(\mathrm{Sp}(u^{-1}) \subset \{0\}\) (prop. 15, a)), we have \(\mathrm{Sp}(u^{-1}) = \{0\}\) (I, p. 26, cor. 1). Since \(u^{-1}\) is normal, this implies that \(u^{-1} = 0\) (I, p. 110, example 1), which is a contradiction since E is not zero.

To prove \(b\) ), one may assume that \(u\) is the multiplication operator \(m_g\) on \(\mathrm{L}^2 (\mathrm{X},\mu)\) , where \(g\) is a continuous function on a locally compact topological space X equipped with a positive measure \(\mu\) (th. 1 of IV, p. 266). Let \(\lambda \in \mathbf{C} - \mathrm{Sp}(u)\) and let \(\delta >0\) be the distance from \(\lambda\) to the spectrum of \(u\) . To prove that \(\| \mathrm{R}(u,\lambda)\| = \delta^{-1}\) , we reduce to the case \(\lambda = 0\) by replacing \(u\) by \(u - \lambda 1_{\mathrm{E}}\) . The real number \(\delta\) is then the distance from 0 to the spectrum of \(m_g\) . Since the latter coincides with the \(\mu\)-essential image of \(g\) , the result is a consequence of the lemma 3 of IV, p. 182.

Finally, the assertion \(c\) ) follows from \(b\) ) and the definition of \(\mathrm{PSp}_{\varepsilon}(u)\) (IV, p. 250, def. 9).

COROLLARY 3. --- Let u be a partial normal operator on a Hilbert space E. Let \(f \in \mathcal{L}_{\mathrm{u}}(\mathrm{Sp}(u))\) . For all x and y in E, the spectral measure of \((x, y)\) relative to \(f(u)\) is the image measure \(f(\mu)\) , where \(\mu\) is the spectral measure of \((x, y)\) relative to u.

In view of the spectral theorem (th. 1 of IV, p. 266), one may suppose that \(u\) is the multiplication operator \(m_g\) on \(\mathrm{L}^2 (\mathrm{X},\mu)\) , where X is a locally compact topological space, \(\mu\) is a positive measure on X and \(g\in \mathcal{C}(\mathrm{X})\) . Let \(f_{1}\) and \(f_{2}\) be in \(\mathcal{L}^2 (\mathrm{X},\mu)\) with classes \(\widetilde{f}_1\) and \(\widetilde{f}_2\) in \(\mathrm{L}^2 (\mathrm{X},\mu)\) . Since \(f(m_g) = m_{f\circ g}\) (prop. 4, c)), the spectral measure of \((\widetilde{f}_1,\widetilde{f}_2)\) relative to \(f(m_g)\) is the image measure \((f\circ g)(\overline{f}_1f_2\cdot \mu)\) (lemma 4, b) of IV, p. 269). This measure is equal to the image measure \(f(g(\overline{f}_1f_2\cdot \mu))\) (INT, V, p. 72, § 6, n° 4, prop. 4, a)), whence the assertion (lemma 4, b) of IV, p. 269).

COROLLARY 4. --- Let \(g \in \mathcal{L}_{\mathrm{u}}(\mathrm{Sp}(f(u)))\) . We have \(g(f(u)) = (g \circ f)(u)\) .

One has \(g \circ f \in \mathcal{L}_{\mathrm{u}}(\mathrm{Sp}(u))\) (lemma 5 of IV, p. 184). For all \(x\) and \(y\) in E, denote by \(\mu_{x,y}^{\prime}\) the spectral measure of \((x,y)\) relative to \(f(u)\) . By the preceding corollary and INT, V, p. 71, § 6, n°2, th. 1, one has \(g \in \mathcal{L}^{2}(\mathrm{Sp}(f(u)), \mu_{x,y}^{\prime})\) if and only if \(g \circ f \in \mathcal{L}^{2}(\mathrm{Sp}(u), \mu_{x,y})\) , so the domain of \(g(f(u))\) is equal to the domain of \((g \circ f)(u)\) . For all \(x \in \mathrm{E}\) and \(y \in \operatorname{dom}(g(f(u)))\) , one then has the formula

\[
\begin{array}{c} \langle x \mid g (f (u)) y \rangle = \int_ {\mathrm{Sp} (f (u))} g \mu_ {x, y} ^ {\prime} \\ = \int_ {\mathrm{Sp} (f (u))} g f (\mu_ {x, y}) = \langle x \mid (g \circ f) y \rangle \end{array}
\]

(loc. cit.), whence \(g(f(u)) = (g \circ f)(u)\) .

PROPOSITION 5. --- Let u be a normal partial operator on a Hilbert space E.

\begin{enumerate}
\def\labelenumi{\alph{enumi})}
\item
  \(If f \in \mathcal{L}_{\mathrm{u}}^{\infty}(\mathrm{Sp}(u)),\) then \(f(u) \in \mathcal{L}(\mathrm{E})\) ;
\item
  \(If f \in \mathcal{L}_{\mathrm{u}}^{\infty}(\mathrm{Sp}(u))\) and \(g \in \mathcal{L}_{\mathrm{u}}(\mathrm{Sp}(u))\) , then \(f(u) \circ g(u) \subset (fg)(u)\) ;
\item
  The map \(f \mapsto f(u)\) from \(\mathcal{L}_{\mathrm{u}}^{\infty}(\mathrm{Sp}(u))\) into \(\mathcal{L}(\mathrm{E})\) is a continuous unital morphism of star algebras. In particular, one has \(\|f(u)\|\leqslant\|f\|_{\infty}\) for \(f\in\mathcal{L}_{\mathrm{u}}^{\infty}(\mathrm{Sp}(u))\) ;
\item
  If \(u \in \mathcal{L}(\mathrm{E})\) , then for all \(f \in \mathcal{L}_{\mathrm{u}}^{\infty}(\mathrm{Sp}(u))\) , the endomorphism \(f(u)\) belongs to the bicommutant of \(u\) in \(\mathcal{L}(\mathrm{E})\) .
\end{enumerate}

Let \(f \in \mathcal{L}_{\mathrm{u}}^{\infty}(\mathrm{Sp}(u))\) . We have \(\mathrm{D}_f = \mathrm{E}\) and \(f(u)\) is a continuous endomorphism of \(\mathrm{E}\) by formula (8), p. 272, whence the assertion \(a\) ).

Let \(g \in \mathcal{L}_{\mathrm{u}}(\mathrm{Sp}(u))\) Let \(y \in \mathrm{D}_g\) . We have \(y \in \mathrm{D}_{fg}\) and, for all \(x \in \mathrm{E}\) , it follows that \(\langle \overline{f}(u)x | g(u)y \rangle = \langle x | (fg)(u)y \rangle\) by formula (7), p. 271, whence \(f(u)(g(u)y) = (fg)(u)(y)\) , which proves \(b\) .

By \(a\) ) and \(b\) ), the map \(f \mapsto f(u)\) of \(\mathcal{L}_{\mathrm{u}}^{\infty}(\mathrm{Sp}(u))\) in \(\mathcal{L}(\mathrm{E})\) is a unital morphism of involutive algebras; consequently, it is a continuous morphism of norm \(\leqslant 1\) (I, p. 104, prop. 2), whence the assertion \(c\) ).

Suppose that \(u\) is an endomorphism of E. Let \(v \in \mathcal{L}(\mathrm{E})\) commuting with \(u\) . We then have \(v \circ f(u) = f(u) \circ v\) for \(f \in \mathcal{C}(\mathrm{Sp}(u))\) by the properties of the continuous functional calculus (I, p. 110, remark). Let \(x\) and \(y\) be in E. The formulas

\[
\langle x \mid (v \circ f (u)) y \rangle = \langle x \mid (f (u) \circ v) y \rangle ,\tag{9}
\]

valid for every function \(f \in \mathcal{C}(\mathrm{Sp}(u))\) , mean that the spectral measures of \((v^{*}(x), y)\) and of \((x, v(y))\) relative to u are equal. This equality implies, by formula (6), p. 271, that formula (9) is valid for all \(f \in \mathcal{L}_{\mathrm{u}}^{\infty}(\mathrm{Sp}(u))\) . We therefore have \(v \circ f(u) = f(u) \circ v\) .

COROLLARY. --- Let \(f\) and \(g\) in \(\mathcal{L}_{\mathrm{u}}(\mathrm{Sp}(u))\) and \((x,y)\in\mathrm{D}_{f}\times\mathrm{D}_{g}\) . The spectral measure of \((f(u)x,g(u)y)\) relative to \(u\) is the measure \(\overline{f}g\cdot\mu_{x,y}\) , where \(\mu_{x,y}\) is the spectral measure of \((x,y)\) relative to \(u\) .

Denote by \(\nu\) the spectral measure of \((f(u)x,g(u)y)\) relative to \(u\) . For every function \(\varphi \in \mathcal{H}(\mathrm{Sp}(u))\) , we have

\[
\begin{array}{r l} & {\int_ {\mathrm{Sp} (u)} \varphi \nu = \langle f (u) x | \varphi (u) (g (u) y) \rangle} \\ & {\qquad = \langle f (u) x | (\varphi g) (u) y \rangle = \int_ {\mathrm{Sp} (u)} \overline {{f}} \varphi g \mu_ {x, y},} \end{array}
\]

(prop. 5, b)) whence \(\nu = \overline{f} g \cdot \mu_{x,y}\) .

PROPOSITION 6. --- Let u be a normal partial operator on a Hilbert space E. Let \((f_n)_{n \in \mathbf{N}}\) be a sequence in \(\mathcal{L}_{\mathrm{u}}(\mathrm{Sp}(u))\) which converges pointwise to \(f \in \mathcal{L}_{\mathrm{u}}(\mathrm{Sp}(u))\) and such that there exists \(g \in \mathcal{L}_{\mathrm{u}}(\mathrm{Sp}(u))\) satisfying \(|f_n| \leqslant g\) for all \(n \in \mathbf{N}\) . We then have \(\operatorname{dom}(g(u)) \subset \operatorname{dom}(f_n(u))\) for all \(n \in \mathbf{N}\) and \(\operatorname{dom}(g(u)) \subset \operatorname{dom}(f(u))\) . Moreover, for every element y of the domain of \(g(u)\) , we have

\[
f (u) (y) = \lim _ {n \to + \infty} f _ {n} (u) (y).
\]

In particular, if \(f_{n} \in \mathcal{L}_{\mathrm{u}}^{\infty}(\mathrm{Sp}(u))\) for all \(n \in \mathbf{N}\) and if the functions \(f_{n}\) are uniformly bounded, then \(f_{n}(u)\) converges to \(f(u)\) in the space \(\mathcal{L}(\mathrm{E})\) endowed with the topology of pointwise convergence.

Denote by \(\mu_{x,y}\) (resp. \(\mu_{x}\) ) the spectral measure of \((x,y)\) (resp. of x) relative to u.

Let \(y \in \mathrm{dom}(g(u))\) , so that \(g \in \mathcal{L}^2(\mathrm{Sp}(u), \mu_y)\) . The condition \(|f_n| \leqslant g\) implies that \(f_n \in \mathcal{L}^2(\mathrm{Sp}(u), \mu_y)\) , hence \(y \in \mathrm{dom}(f_n(u))\) .

Since \((f_{n})\) converges simply to f and that \(|f_{n}| \leqslant g\) , we have \(f \in \mathcal{L}^{2}(\mathrm{Sp}(u), \mu_{y})\) by the Lebesgue theorem (INT, IV, p. 137, § 3, n° 7, th. 6), hence \(y \in \mathrm{dom}(f(u))\) . Moreover, the sequence \((f_{n})\) converges to f in \(\mathcal{L}^{2}(\mathrm{Sp}(u), \mu_{y})\) , so the norm of \(f_{n}(u)y\) converges to the norm of \(f(u)y\) .

Let \(x \in \mathrm{E}\) . The functions \(f\) and \(g\) , as well as the functions \(f_n\) for all \(n \in \mathbf{N}\) , belong to \(\mathcal{L}^1(\mathrm{Sp}(u), \mu_{x,y})\) (prop. 3). By the Lebesgue theorem (INT, IV, loc. cit.), the sequence \((f_n)\) converges to \(f\) in \(\mathcal{L}^1(\mathrm{Sp}(u), \mu_{x,y})\) , whence

\[
\langle x \mid f _ {n} (u) y \rangle = \int_ {\operatorname{Sp} (u)} f _ {n} \mu_ {x, y} \rightarrow \int_ {\operatorname{Sp} (u)} f \mu_ {x, y} = \langle x \mid f (u) y \rangle .
\]

One concludes that \(f_{n}(u)(y)\) converges to \(f(u)(y)\) (EVT, V, p. 17, prop. 10).

PROPOSITION 7. --- Let X be a locally compact topological space and let ν be a measure on X. Let g: C × X → C be a continuous function with compact support. For z ∈ C, set

\[
h (z) = \int_ {\mathrm{X}} g (z, x) d \nu (x).
\]

\begin{enumerate}
\def\labelenumi{\alph{enumi})}
\item
  The map h from C to C is continuous and bounded;
\item
  The map from X to \(\mathcal{L}(\mathrm{E})\) defined by \(x\mapsto g(u,x)\) is \(\nu\) -integrable and we have
\end{enumerate}

\[
h (u) = \int_ {\mathrm{X}} g (u, x) d \nu (x).
\]

The function h is bounded because g is continuous with compact support, and it is continuous by INT, IV, p. 144, § 4, n° 3, cor. 1. Since g is continuous and has compact support, the map \(x \mapsto g(u, x)\) is continuous from X into \(\mathcal{L}(\mathrm{E})\) . This map has compact support, hence is bounded and integrable on X with respect to \(\nu\) . Denote by

\[
v = \int_ {\mathrm{X}} g (u, x) d \nu (x) \in \mathscr {L} (\mathrm{E}).
\]

Let y and z be elements of E, and let \(\mu\) be the spectral measure of \((y, z)\) relative to u. We have

\[
\langle y \mid v (z) \rangle = \int_ {\mathrm{X}} \langle y \mid g (u, x) z \rangle d \nu (x) = \int_ {\mathrm{X}} \left(\int_ {\operatorname{Sp} (u)} g (\lambda , x) d \mu (\lambda)\right) d \nu (x)
\]

(formula (6), p. 271). Since \(g \in \mathcal{K}(\mathbf{C} \times \mathbf{X})\) , it follows that

\[
\langle y \mid v (z) \rangle = \int_ {\operatorname{Sp} (u)} \Bigl (\int_ {\mathrm{X}} g (\lambda , x) d \nu (x) \Bigr) d \mu (\lambda) = \langle y \mid h (u) z \rangle
\]

by INT, III, p. 84, § 4, n° 1, th. 2 and formula (6), p. 271. This proves that \(v = h(u)\) , as desired.

